\section{\label{sec:introduction} Introduction}

Spinors were discovered by {\'E}lie Cartan \cite{Cartan1913groupes}
while he was working on the linear representation of simple groups.
They are the most basic representation of the
rotation and Lorentz groups and all other representations can be obtained
as tensor products of spinors \cite{cartan2012theory}.
For this reason, they are often loosely described as the square
roots of Euclidean (or Minkowskian) vectors \cite{farmelo2009strangest}.
Spinors are difficult to picture not only because
their components are in general complex, but also because when the Euclidean space
is rotated by an angle of $2\pi$, they do not return to their initial state (like vectors do),
but acquire a multiplicative phase of $-1$. Only rotations
of multiples of $4\pi$ are thus equivalent to the identity operator on spinors.

The standard approach for visualizing the state of a two-level quantum system is the Bloch sphere representation. While powerful and widely used, the Bloch sphere has a fundamental limitation: it represents only the \emph{direction} associated with a spinor (via the Bloch vector $\vv$), collapsing the global phase information. Furthermore, the Bloch sphere does not make explicit how the individual spinor components contribute to the overall state.

Spinors are not only intimately related to several group-theoretical
and geometrical objects, such as quaternions, Hopf fibrations,
stereographic projections, and M\"obius transformations
\cite{hanson2006visualizing,penrose1984spinors,needham2021visual},
but also of great importance for physics, as they are needed
to correctly describe the wavefunction of $\half$ particles
\cite{hladik1996spineurs,schwichtenberg2015physics}.

In this work, we develop a geometric representation of
two-dimensional spinors that overcomes these limitations.
Starting from simple considerations about vectors, reference
frames, and half-angle constructions, we build a visualization
where:
\begin{itemize}
    \item The \emph{moduli} of the spinor components are encoded as the radii of two circle sections on a sphere;
    \item The \emph{relative phase} between components appears as a geometric rotation of these circles;
    \item The \emph{global phase} is also geometrically represented;
    \item The characteristic $4\pi$-periodicity of spinors is made visually manifest.
\end{itemize}

The paper is organized as follows. In Section~\ref{sec:planar} we introduce a planar half-angle construction that serves as the foundation for the three-dimensional visualization. Section~\ref{sec:hyperchords} promotes this construction to $\Rt$, defining the ``hyperchord'' circles that form our spinor representation. Section~\ref{sec:su2_viz} connects the geometric construction to the standard spinor formalism and shows how $\sut$ rotations act on the visualization. Section~\ref{sec:algorithm} presents the rendering algorithm with pseudocode. We conclude with a discussion in Section~\ref{sec:discussion}.
An open-source Python implementation of the visualization, together with interactive Jupyter notebooks, is available at \url{https://github.com/al-rigazzi/spinor_viz}~\cite{spinor_viz_github}.

The reader interested in a comprehensive exposure of spinors
and their algebra is referred to articles \cite{steane2013introduction,wheler2000spinor}
and monographs \cite{lounesto2001clifford,penrose1984spinors,hall2015lie} on the subject.
